% !TeX root = ../main.tex
% Appendix A4: Quantum Trading Strategy
% Batch #4.9b — Notation hygiene, aligned with Fig 2 (6.7d lead) and Tab 6 (marginal case study)

\section{Quantum Trading Strategy}\label{app:trading}

A trading strategy based on QST exploits the leading indicator properties of
$\mathcal{C}(\rho)$ and $S(\rho)$ established in
Sections~\ref{M-sec:coherence-2023} and~\ref{M-sec:entropy-2023} to generate equity
option signals.

\subsection{Signal Construction}

The trading signal is constructed from the daily QST estimates of the market
state. Let $\theta_t$ be the distress amplitude, $\phi_t$ the contagion phase,
and $\hbar_{\mathrm{econ}}$ the economic Planck constant. We define three trading
signals:

\begin{enumerate}
    \item \textbf{Coherence signal:}
        $S_t^{\mathrm{coh}} = \mathcal{C}(\rho_t) - \overline{\mathcal{C}}$,
        where $\overline{\mathcal{C}}$ is the rolling 30-day mean of the quantum
        coherence. A positive signal indicates increasing systemic contagion.
    \item \textbf{Entropy signal:} $S_t^{\mathrm{ent}} = S(\rho_t) - \bar{S}$,
        where $\bar{S}$ is the rolling 30-day mean of the von Neumann entropy.
        A positive signal indicates increasing structural uncertainty.
    \item \textbf{Combined signal:}
        $S_t^{\mathrm{comb}} = w_1 S_t^{\mathrm{coh}} + w_2 S_t^{\mathrm{ent}}$,
        where $w_1$ and $w_2$ are weights estimated from the in-sample period.
\end{enumerate}

\subsection{Portfolio Construction}

The trading strategy is implemented as a long--short portfolio of equity options.
When the combined signal exceeds a threshold $\tau_{\mathrm{sig}}$, the strategy
takes a long position in out-of-the-money put options on the KRE Regional Banking
ETF and a short position in out-of-the-money put options on the S\&P~500 Index.
This pair trade isolates the regional banking sector risk captured by the
quantum state.

\subsection{Empirical Signal Timing}

For the 2023 US regional banking crisis case study
(Section~\ref{M-sec:coherence-2023}), the coherence signal $S_t^{\mathrm{coh}}$
would generate an entry indication on March~9, 2023 (one trading day before SVB's
collapse on March~10). Averaged across the three failure events
(SVB March~10, Signature Bank March~12, First Republic Bank May~1), the
coherence signal leads the corresponding classical Pearson-correlation crossing
by 7~trading days. The entropy signal $S_t^{\mathrm{ent}}$ provides an even
earlier confirmation, crossing its 95th-percentile baseline threshold on
February~17, 2023---15~trading days before the SVB event.

\subsection{Case-Study Performance and Caveats}

We evaluate the signal's association with option-market P\&L for the KRE
regional-banking option chain during the crisis window. Because
2023-03-10-vintage KRE option chain data are not publicly available (see
Section~S.2 on the canonical smile reconstruction), a
full-cross-strike hedging back-test is beyond the scope of the current
reproducibility package. The single-path hedging case study in
Online Supplement~S.2 (Table~\ref{tab:hedging-rmshe}) demonstrates that a signal-triggered quantum
optimal hedge (per Proposition~\ref{M-prop:quantum-optimal-hedge}) yields a
marginal RMSHE reduction relative to the BSM baseline over the 20-day stress
window (2023-02-27 to 2023-03-24), with the correction concentrated in the
post-SVB OTM segment where $x = \ln(S_t/K)$ is non-zero and $\tau$ is short.

\emph{Disclaimer:} The signal-timing results above are based on a single
crisis episode and should be interpreted as an illustrative case study rather
than an out-of-sample validation. A comprehensive out-of-sample evaluation
across multiple crisis and non-crisis periods, using historical option-chain
data across the equity-financial sector, is left for future work.
